\documentclass[12pt,a4paper]{article}
\usepackage[T2A]{fontenc}
\usepackage[utf8]{inputenc}
\usepackage[russian]{babel}
\usepackage{amsmath,amssymb}
\usepackage{array,booktabs,multirow,longtable,tabularx}
\usepackage[left=2cm,right=2cm,top=2cm,bottom=2cm]{geometry}
\usepackage{graphicx}
\usepackage{float}
\usepackage{caption}
\usepackage{tikz}
\usepackage{pgfplots}
\usepackage{hyperref}
\pgfplotsset{compat=1.18}
\usepackage{indentfirst}
\renewcommand{\arraystretch}{1.18}
\setlength{\tabcolsep}{6pt}
\hypersetup{colorlinks=true,urlcolor=blue,linkcolor=black}

% === ПОДКЛЮЧЕНИЕ LISTINGS ДЛЯ КОДА ===
\usepackage{listings}
\usepackage{xcolor}

% Стиль для Python с поддержкой кириллицы в комментариях
\lstdefinestyle{pythonstyle}{
    language=Python,
    basicstyle=\ttfamily\small,
    keywordstyle=\color{blue}\bfseries,
    commentstyle=\color{green!60!black},
    stringstyle=\color{red!80!black},
    numbers=left,
    numberstyle=\tiny\color{gray},
    stepnumber=1,
    numbersep=8pt,
    breaklines=true,
    breakatwhitespace=true,
    showstringspaces=false,
    tabsize=4,
    frame=single,
    rulecolor=\color{black!30},
    backgroundcolor=\color{gray!5},
    captionpos=b,
    keepspaces=true,
    inputencoding=utf8,
    extendedchars=true,
    literate=
        {а}{{\selectfont\char224}}1{б}{{\selectfont\char225}}1{в}{{\selectfont\char226}}1
        {г}{{\selectfont\char227}}1{д}{{\selectfont\char228}}1{е}{{\selectfont\char229}}1
        {ё}{{\"e}}1{ж}{{\selectfont\char230}}1{з}{{\selectfont\char231}}1
        {и}{{\selectfont\char232}}1{й}{{\selectfont\char233}}1{к}{{\selectfont\char234}}1
        {л}{{\selectfont\char235}}1{м}{{\selectfont\char236}}1{н}{{\selectfont\char237}}1
        {о}{{\selectfont\char238}}1{п}{{\selectfont\char239}}1{р}{{\selectfont\char240}}1
        {с}{{\selectfont\char241}}1{т}{{\selectfont\char242}}1{у}{{\selectfont\char243}}1
        {ф}{{\selectfont\char244}}1{х}{{\selectfont\char245}}1{ц}{{\selectfont\char246}}1
        {ч}{{\selectfont\char247}}1{ш}{{\selectfont\char248}}1{щ}{{\selectfont\char249}}1
        {ъ}{{\selectfont\char250}}1{ы}{{\selectfont\char251}}1{ь}{{\selectfont\char252}}1
        {э}{{\selectfont\char253}}1{ю}{{\selectfont\char254}}1{я}{{\selectfont\char255}}1
        {А}{{\selectfont\char192}}1{Б}{{\selectfont\char193}}1{В}{{\selectfont\char194}}1
        {Г}{{\selectfont\char195}}1{Д}{{\selectfont\char196}}1{Е}{{\selectfont\char197}}1
        {Ё}{{\"E}}1{Ж}{{\selectfont\char198}}1{З}{{\selectfont\char199}}1
        {И}{{\selectfont\char200}}1{Й}{{\selectfont\char201}}1{К}{{\selectfont\char202}}1
        {Л}{{\selectfont\char203}}1{М}{{\selectfont\char204}}1{Н}{{\selectfont\char205}}1
        {О}{{\selectfont\char206}}1{П}{{\selectfont\char207}}1{Р}{{\selectfont\char208}}1
        {С}{{\selectfont\char209}}1{Т}{{\selectfont\char210}}1{У}{{\selectfont\char211}}1
        {Ф}{{\selectfont\char212}}1{Х}{{\selectfont\char213}}1{Ц}{{\selectfont\char214}}1
        {Ч}{{\selectfont\char215}}1{Ш}{{\selectfont\char216}}1{Щ}{{\selectfont\char217}}1
        {Ъ}{{\selectfont\char218}}1{Ы}{{\selectfont\char219}}1{Ь}{{\selectfont\char220}}1
        {Э}{{\selectfont\char221}}1{Ю}{{\selectfont\char222}}1{Я}{{\selectfont\char223}}1
}
\lstset{style=pythonstyle}
% =====================================

\begin{document}

\noindent
\begin{minipage}[t]{0.6\textwidth}
    \vspace{0pt}
    Университет ИТМО \\
    Факультет «Безопасность информационных технологий» (ФБИТ)
\end{minipage}%
\hfill%
\begin{minipage}[t]{0.3\textwidth}
    \vspace{-14pt}
    \centering
    \includegraphics[width=4cm]{itmo.png}
\end{minipage}
\rule{\textwidth}{1.5pt}
\vspace{0.01cm}

\noindent
Группа \underline{\hspace{3cm} N3145 \hspace{1.5cm}} \hfill К работе допущены \underline{\hspace{4cm}}\\
Студент \underline{\hspace{0.2cm} Савченко А.Е. \hspace{0.2cm}} \hfill Работа выполнена \underline{\hspace{3.4cm}}\\
Преподаватель \underline{\hspace{0.935cm} Замолоцких С.С. \hspace{1.5cm}} \hfill Отчет принят \underline{\hspace{4.84cm}}\\

\vspace{0.01cm}

\begin{center}
    \Large\textbf{Рабочий протокол и отчет} \\
    \Large\textbf{по проектной работе по физике}
\end{center}
\noindent\rule{\textwidth}{1.5pt}%
\vspace{-0.3cm}
\begin{center}
Визуализация относительности движения на прямой
\end{center}
\vspace{-0.3cm}
\noindent\rule{\textwidth}{1.5pt}

\section{Цель работы}
Разработка интерактивной компьютерной модели, демонстрирующей относительность положения и скорости двух тел при равномерном движении вдоль одной прямой.

\section{Задачи}
\begin{enumerate}
    \item Реализовать математическую модель движения двух машин по прямому шоссе при постоянных скоростях.
    \item Обеспечить возможность задания начальных координат и скоростей пользователем.
    \item Реализовать выбор системы отсчета: дорога, машина A, машина B.
    \item Выполнить пересчет координат и скоростей при переходе между системами отсчета.
    \item Визуализировать движение и численные параметры
\end{enumerate}

\section{Объект исследования}
Классическое равномерное прямолинейное движение двух тел и преобразование координат и скоростей при переходе между инерциальными системами отсчета в рамках преобразований Галилея.

\section{Метод исследования}
В работе использован метод математического моделирования. Движение каждой машины описывается уравнением равномерного прямолинейного движения. Далее выполняется перенос начала координат в выбранную систему отсчета и пересчет скоростей относительно нового наблюдателя. Полученные зависимости выводятся в численной форме, отображаются на анимированной сцене. $x(t)$, $x'(t)$, $v(t)$ и $\Delta x(t)$.

\section{Рабочие формулы и исходные соотношения}
Для лабораторной системы отсчета, связанной с дорогой, координаты машин задаются формулами
\begin{equation}
 x_A(t)=x_{0A}+v_A t,
\end{equation}
\begin{equation}
 x_B(t)=x_{0B}+v_B t.
\end{equation}

Расстояние между машинами вдоль дороги:
\begin{equation}
 \Delta x(t)=x_B(t)-x_A(t).
\end{equation}

Относительная скорость машины A относительно машины B:
\begin{equation}
 v_{\text{отн}}=v_A-v_B.
\end{equation}

При переходе в систему отсчета, движущуюся вместе с опорным телом, используется галилеево преобразование
\begin{equation}
 x'=x-x_{\text{опор}}(t),
\end{equation}
\begin{equation}
 v'=v-v_{\text{опор}}.
\end{equation}

Если система отсчета связана с машиной A, то
\begin{equation}
 x'_A=0,
\end{equation}
\begin{equation}
 x'_B=x_B-x_A,
\end{equation}
\begin{equation}
 v'_A=0,
\end{equation}
\begin{equation}
 v'_B=v_B-v_A.
\end{equation}

Если система отсчета связана с машиной B, то
\begin{equation}
 x'_B=0,
\end{equation}
\begin{equation}
 x'_A=x_A-x_B,
\end{equation}
\begin{equation}
 v'_B=0,
\end{equation}
\begin{equation}
 v'_A=v_A-v_B.
\end{equation}

Для рассматриваемой модели время во всех выбранных системах отсчета одинаково:
\begin{equation}
 t'=t.
\end{equation}

Таким образом, физический смысл проекта состоит в том, что движение и покой являются относительными: одна и та же картина движения описывается по-разному в зависимости от выбора наблюдателя.

\section{Педагогическая ценность модели}
Разработанный продукт полезен не только как программный проект, но и как учебный инструмент. Его ценность состоит в следующем:
\begin{enumerate}
    \item Пользователь может \textbf{наглядно увидеть}, что движение не является абсолютным.
    \item Формулы кинематики сразу связываются с \textbf{визуальным результатом} на экране.
    \item Переключение между системами отсчета позволяет \textbf{физически почувствовать} смысл выражений $x'=x-x_{\text{опор}}(t)$ и $v'=v-v_{\text{опор}}$.
    \item Проект можно использовать как основу для дальнейшего развития: добавить ускорение, большее число тел, сравнение с другими типами движения.
\end{enumerate}

\section{Окончательные результаты}
В результате выполнения проектной работы создан интерактивный продукт, который:
\[
\text{моделирует } x_A(t),\ x_B(t),\ \Delta x(t),\ v_{\text{отн}},\ x'(t),\ v'(t)
\]
и демонстрирует относительность движения при выборе различных систем отсчета.

Основной содержательный результат работы состоит в подтверждении следующего вывода:
\[
\text{движение и покой тел зависят от выбранной системы отсчета.}
\]

\section{Выводы и анализ результатов работы}
В ходе выполнения проектной работы была разработана и проанализирована интерактивная модель движения двух машин по прямому шоссе. В модели реализованы все основные требования задачи: ввод начальных координат и скоростей, визуализация движения, пересчет параметров при выборе новой системы отсчета, вывод численных характеристик.

Физическая часть проекта основана на классической кинематике и преобразованиях Галилея. Проведенные расчеты показывают, что при смене наблюдателя изменяются координаты и скорости отдельных тел, но сохраняется физический смысл их взаимного движения. Именно это и является главным содержанием темы «относительность движения на прямой».

Практическая значимость работы состоит в том, что полученный продукт позволяет изучать тему не только по формулам, но и через непосредственное взаимодействие с моделью. Пользователь может менять параметры, запускать сценарии и сразу видеть результат на анимированной сцене и графиках. Благодаря этому проект выступает как полноценное цифровое учебное пособие по теме.

\section{Листинг кода}
\label{sec:listing}

Ниже приведён полный исходный код программы, реализующей интерактивную визуализацию относительности движения. Программа написана на языке Python с использованием библиотек \texttt{matplotlib} и \texttt{numpy}.

\begin{lstlisting}[caption={Основной код симуляции (main.py)}, label={lst:main}]
import sys

interactive_backends = ['TkAgg', 'Qt5Agg', 'QtAgg', 'WXAgg', 'GTK3Agg', 'GTK4Agg']
backend_set = False

for backend in interactive_backends:
    try:
        import matplotlib
        matplotlib.use(backend, force=True)
        backend_set = True
        break
    except Exception:
        continue

import numpy as np
import matplotlib.pyplot as plt
from matplotlib.animation import FuncAnimation
from matplotlib.widgets import Button


class RelativeMotionSimulation:
    def __init__(self):
        print("Симуляция относительности движения")
        print("-" * 50 + "\n")

        self.x01 = self.get_float_input("Начальная координата Машины 1 (м): ", 0)
        self.v1 = self.get_float_input("Скорость Машины 1 (м/с): ", 10)
        self.x02 = self.get_float_input("Начальная координата Машины 2 (м): ", 100)
        self.v2 = self.get_float_input("Скорость Машины 2 (м/с): ", 20)
        self.duration = self.get_float_input("Длительность симуляции (с): ", 10)

        print(f"\n Запуск визуализации\n")

        self.fps = 30
        self.total_frames = int(self.duration * self.fps)
        self.current_so = 'road'
        self.trail1 = []
        self.trail2 = []
        self.margin = 100

        self.fig, self.ax = plt.subplots(figsize=(12, 6))
        plt.subplots_adjust(bottom=0.25)

        self.road_line, = self.ax.plot([], [], 'k--', linewidth=1, alpha=0.5)
        self.car1_point, = self.ax.plot([], [], 'ro', markersize=15, label='Машина 1', zorder=5)
        self.car2_point, = self.ax.plot([], [], 'bo', markersize=15, label='Машина 2', zorder=5)
        self.trail1_line, = self.ax.plot([], [], 'r-', alpha=0.4, linewidth=2)
        self.trail2_line, = self.ax.plot([], [], 'b-', alpha=0.4, linewidth=2)

        self.info_text = self.ax.text(0.02, 0.95, '', transform=self.ax.transAxes,
                                      verticalalignment='top', fontsize=9,
                                      bbox=dict(boxstyle='round', facecolor='wheat', alpha=0.7))

        self.ax.set_ylim(-0.5, 0.5)
        self.ax.set_yticks([])
        self.ax.set_xlabel('Координата (м)', fontsize=12)
        self.ax.set_title('Относительность движения', fontsize=14, fontweight='bold')
        self.ax.grid(True, axis='x', alpha=0.3)
        self.ax.legend(loc='upper right', fontsize=10)

        # Кнопки переключения СО
        self.ax_btn_road = plt.axes([0.1, 0.05, 0.22, 0.075])
        self.ax_btn_car1 = plt.axes([0.39, 0.05, 0.22, 0.075])
        self.ax_btn_car2 = plt.axes([0.68, 0.05, 0.22, 0.075])

        self.btn_road = Button(self.ax_btn_road, 'СО: Дорога', color='lightgray', hovercolor='lightgreen')
        self.btn_car1 = Button(self.ax_btn_car1, 'СО: Машина 1', color='lightcoral', hovercolor='red')
        self.btn_car2 = Button(self.ax_btn_car2, 'СО: Машина 2', color='lightskyblue', hovercolor='blue')

        self.btn_road.on_clicked(self.set_so_road)
        self.btn_car1.on_clicked(self.set_so_car1)
        self.btn_car2.on_clicked(self.set_so_car2)

        # Кнопки масштаба
        self.ax_btn_zoom_in = plt.axes([0.88, 0.88, 0.05, 0.04])
        self.ax_btn_zoom_out = plt.axes([0.94, 0.88, 0.05, 0.04])

        self.btn_zoom_in = Button(self.ax_btn_zoom_in, '+', color='lightgreen', hovercolor='green')
        self.btn_zoom_out = Button(self.ax_btn_zoom_out, '-', color='lightcoral', hovercolor='red')

        self.btn_zoom_in.on_clicked(self.zoom_in)
        self.btn_zoom_out.on_clicked(self.zoom_out)

    def get_float_input(self, prompt, default):
        while True:
            try:
                user_input = input(prompt).strip()
                if user_input == '':
                    return default
                return float(user_input)
            except ValueError:
                print("Введите число или нажмите Enter.")

    def get_positions_road_so(self, t):
        return self.x01 + self.v1 * t, self.x02 + self.v2 * t

    def get_positions_relative(self, t, so):
        x1_road, x2_road = self.get_positions_road_so(t)
        if so == 'road':
            return x1_road, x2_road
        elif so == 'car1':
            return 0, x2_road - x1_road
        elif so == 'car2':
            return x1_road - x2_road, 0

    def update(self, frame):
        t = frame / self.fps
        x1, x2 = self.get_positions_relative(t, self.current_so)

        self.trail1.append(x1)
        self.trail2.append(x2)

        self.car1_point.set_data([x1], [0])
        self.car2_point.set_data([x2], [0])
        self.trail1_line.set_data(self.trail1, [0] * len(self.trail1))
        self.trail2_line.set_data(self.trail2, [0] * len(self.trail2))

        if self.current_so == 'road':
            center = (x1 + x2) / 2
            self.ax.set_xlim(center - self.margin, center + self.margin)
        else:
            self.ax.set_xlim(-self.margin, self.margin)

        self.road_line.set_data(self.ax.get_xlim(), [0, 0])

        x1_road, x2_road = self.get_positions_road_so(t)
        v_rel = self.v2 - self.v1 if self.current_so != 'car2' else self.v1 - self.v2

        info = (f"Время: {t:.1f} с | СО: {self.get_so_name()} | Масштаб: +/-{self.margin}м\n"
                f"-- В СО Дороги --\n"
                f"M1: x={x1_road:6.1f} м, v={self.v1:5.1f} м/с\n"
                f"M2: x={x2_road:6.1f} м, v={self.v2:5.1f} м/с\n"
                f"-- В текущей СО --\n"
                f"M1: x'={x1:6.1f} м | M2: x'={x2:6.1f} м\n"
                f" V_otn: {v_rel:5.1f} м/с")
        self.info_text.set_text(info)

        return self.car1_point, self.car2_point, self.trail1_line, self.trail2_line, self.info_text, self.road_line

    def get_so_name(self):
        return {'road': 'Дорога', 'car1': 'Машина 1', 'car2': 'Машина 2'}.get(self.current_so, '?')

    def set_so_road(self, event):
        self.current_so = 'road'
        self.reset_trails()
        self.ax.set_title('СО: Дорога (Неподвижная)', fontsize=14, fontweight='bold')
        self.fig.canvas.draw_idle()

    def set_so_car1(self, event):
        self.current_so = 'car1'
        self.reset_trails()
        self.ax.set_title('СО: Машина 1 (Наблюдатель в М1)', fontsize=14, fontweight='bold')
        self.fig.canvas.draw_idle()

    def set_so_car2(self, event):
        self.current_so = 'car2'
        self.reset_trails()
        self.ax.set_title('СО: Машина 2 (Наблюдатель в М2)', fontsize=14, fontweight='bold')
        self.fig.canvas.draw_idle()

    def reset_trails(self):
        self.trail1, self.trail2 = [], []

    def zoom_in(self, event):
        self.margin = max(10, self.margin - 20)
        self.fig.canvas.draw_idle()

    def zoom_out(self, event):
        self.margin += 20
        self.fig.canvas.draw_idle()

    def run(self):
        try:
            self.anim = FuncAnimation(self.fig, self.update, frames=self.total_frames,
                                      interval=1000 / self.fps, blit=True, repeat=False)
            plt.show()
            print("\n Симуляция завершена.")
        except Exception as e:
            print(f"\n Ошибка при запуске: {e}")
            print(" Запускаю режим сохранения кадров...")
            self.save_static_plot()

    def save_static_plot(self):
        try:
            t = self.duration / 2
            x1, x2 = self.get_positions_relative(t, self.current_so)
            self.car1_point.set_data([x1], [0])
            self.car2_point.set_data([x2], [0])
            self.ax.set_title(f'Статичный кадр (t={t:.1f} с)')
            plt.savefig('motion_snapshot.png', dpi=150, bbox_inches='tight')
            print(" Сохранено изображение: motion_snapshot.png")
        except Exception as e:
            print(f" Не удалось сохранить: {e}")


if __name__ == "__main__":
    try:
        sim = RelativeMotionSimulation()
        sim.run()
    except KeyboardInterrupt:
        print("\n Прервано.")
    except Exception as e:
        print(f"\n Критическая ошибка: {e}")
        sys.exit(1)
\end{lstlisting}

\end{document}